\documentclass[11pt]{article}

% --- Packages ---
\usepackage[spanish]{babel}
\usepackage[utf8]{inputenc} % Ensures accents like á, é, ñ work
\usepackage[T1]{fontenc}    % Proper font encoding for European languages
\usepackage[margin=1in]{geometry}
\usepackage{enumitem}
\usepackage{setspace}
\usepackage{titlesec}
\usepackage{tikz}
\usepackage{hyperref}
\hypersetup{
    colorlinks=true,
    linkcolor=blue,
    filecolor=blue,
    urlcolor=blue,
    pdftitle={Overleaf Example},
    pdfpagemode=FullScreen,
}

% --- Formatting ---
\setstretch{1.2} % Better readability
\titleformat{\section}{\normalfont\large\bfseries}{}{0pt}{}
\setlength{\parindent}{0pt}
\setlength{\parskip}{1em}

\begin{document}

% --- Header ---
{\centering
    \LARGE \textbf{KNOWLEDGE TRANSFER AGREEMENT} \\
    \vspace{0.5em}
    \small Este acuerdo es efectivo a partir del \textbf{\today}
\par}

\vspace{2em}

\section{Bases de Datos.}\label{contrato---administraciuxf3n-de-bases-de-datos.}
\vspace{2em}

\noindent Entre los suscritos, por una parte, \textbf{El Profesor}, quien en adelante se denominará el \textbf{El Profesor}, y por la otra, \textbf{El Estudiante} inscrito en la asignatura, quien en adelante se denominará \textbf{El Estudiante}, han convenido celebrar el presente acuerdo que se regirá por las siguientes cláusulas:

\vspace{1em}

\subsection{Horario de Clases:}\label{horario-de-clases}
\begin{itemize}
    \item Las clases están programadas los días Martes y Jueves de 2:00pm a 4:00pm, cada semana durante el segundo semestre académico del 2026.
    \item El inicio de la clase será 15 minutos después de la hora programada con el fin de facilitar el desplazamiento y llegada de todos los estudiantes al aula de clase. Se distribuirán 15 minutos de descanso a lo largo de la clase si hubiese lugar.
    \item En los días asignados para parciales, el inicio y finalización de la clase será en la hora en punto programada (\textbf{2:00pm}) con el fin de que los estudiantes tengan el tiempo reglamentario para el desarrollo de sus actividades evaluativas.
\end{itemize}

\subsection{Asistencia a clase:}\label{asistencia-a-clase}
\begin{itemize}
    \item La clase se desarrolla bajo el principio de responsabilidad del estudiante universitario, por lo tanto no se realiza un control de asistencia a las clases, sin embargo es responsabilidad del estudiante actualizarse de los temas vistos en clase a causa de su inasistencia.
    \item Los laboratorios deben ser defendidos por los miembros del grupo y ambos deben estar presentes.  La defensa es presencial y no se admiten dispositivos electrónicos (celulares, tablets, etc).  La inasistencia a las actividades evaluativas de talleres sin justificación válida será responsabilidad del estudiante asumiendo una penalización de 20\% por cada día de mora en la entrega. No se realizará ningún tipo de actividad de defensa o recuperación de laboratorios después de 7 días posteriores a la fecha inicial de defensa. La evaluación de los talleres será teniendo en cuenta el material entregado en BrightSpace (BS) dentro de las franjas establecidas para dicho fin.
    \item La inasistencia a las actividades evaluativas de parciales será responsabilidad del estudiante asumiendo la calificación mínima posible de acuerdo con la escala de calificación de la universidad, salvo los casos debidamente justificados ante la Dirección de Carrera que permita y avale la realización de un supletorio.
\end{itemize}

\subsection{Plataformas utilizadas para el desarrollo del curso:}\label{plataformas-utilizadas-para-el-desarrollo-del-curso}
\begin{itemize}
    \item Los anuncios y comunicados del curso serán publicados a través del email masivo por medio de BS.
    \item Los archivos entregables (informes, trabajos, tareas, proyectos, etc) también serán gestionados a través de BS.
    \item El material de clase como dispositivas y lecturas serán gestionados a través GitHub (GH) \href{https://github.com/aocalderon/PUJ}{https://github.com/aocalderon/PUJ}.
\end{itemize}

\subsection{Metodología:}\label{metodologuxeda}
\begin{itemize}
    \item Las clases serán realizadas en modalidad \textbf{presencial} de acuerdo con los lineamientos de la universidad.
    \item Para el desarrollo del \textbf{Proyecto Final}, los estudiantes estarán organizados en grupos de trabajo de máximo 4 integrantes.
    \item Todos los archivos calificables \textbf{grupales} deben ser oportunamente cargados en la plataforma de BS por alguno de los integrantes del grupo haciendo uso de los links destinados para tal fin. \textbf{Nota:} La plataforma solo mostrará el ultimo archivo cargado, tengan cuidado con las versiones que suben.
    \item La calificación de talleres y del proyecto \textbf{es de manera grupal}.
    \item La calificación de tareas y parciales \textbf{es de manera individual}.
\end{itemize}

\subsection{Archivos:}\label{archivos}
\begin{itemize}
    \item Las diapositivas presentadas en clase serán cargadas después de finalizar la temática a GH para que el estudiante pueda revisar su contenido posteriormente.
    \item Los archivos de códigos de ejemplo o desarrollo de ejercicios serán cargados a GH para que el estudiante pueda revisar su contenido posteriormente.
    \item El orden de los archivos publicados en GH, estará asociado a la semana académica en el cual se desarrolla su contenido.
    \item Todos los archivos entregables de texto que los estudiantes carguen a la plataforma BS para ser calificados deben estar en formato \textbf{PDF}.
    \item Los archivos especiales como ejercicios de código y similares deberán ser cargados en el formato que indique el profesor.
    \item Todos los archivos deben ser cargados a la plataforma de BS en los \textbf{horarios habilitados} para las entregas. Entregas posteriores \textbf{serán} tenidas en cuenta a con una penalización del 20\% por cada día de mora.
\end{itemize}

\subsection{Desarrollo y entrega de talleres:}\label{desarrollo-y-entrega-de-talleres}
\begin{itemize}
\item El desarrollo de los talleres se realizará durante la sesión de clase  (generalmente).
\item Durante las sesiones de clase los estudiantes, de manera grupal, deberán descargar el archivo del taller y desarrollar los ejercicios propuestos.
\item El link de la asignación para la entrega del informe de laboratorio estará en BS.
\end{itemize}

\subsection{Desarrollo y entrega de Evaluaciones:}\label{desarrollo-y-entrega-de-evaluaciones}
\begin{itemize}
\item Los parciales se desarrollarán durante las jornadas de clase establecidas para tal fin.
\item \textbf{El uso de cualquier tipo de dispositivo móvil o IA durante las evaluaciones está restringido}.
\item Si usted considera que necesita algun tipo de dispositivo o condición particular para el desarrollo de alguna evaluación puede solicitarlo con la debida anticipación a través de programa de inclusión y ajustes razonables del Centro de Identidad y Comunidad.
\end{itemize}

\subsection{Calificación:}\label{calificacion}
\textbf{NOTA}: Las fechas son tentativas.
\begin{itemize}
        \item 20\% → Primer Parcial (Semana 6).
        \item 15\% → Segundo Parcial (Semana 12).
        \item 15\% → Examen Final (Semana 17).
        \item 20\% → Laboratorios (Semanales).
        \item 30\% → Proyecto Final (Semana 11, 3 entregas).
\end{itemize}

\subsection{Contacto y horarios de atención fuera de clase:}\label{contacto-y-horarios-de-atenciuxf3n-fuera-de-clase}
\begin{itemize}
    \item La comunicación Estudiantes-Profesor se hará por medio de correo electrónico: \\  \href{mailto:///andrescalderonr@javeriana.edu.co}{andrescalderonr@javeriana.edu.co} (Andrés Oswaldo Calderón Romero) o por el chat de Teams.
    \item Cualquier correo al profesor debe incluir al inicio del asunto la siguiente combinación de caracteres: \textbf{{[}DBS{]}}.
    \item Todo correo o mensaje de chat será contestado en horario laboral, es decir, de 8:00am a 6:00pm de Lunes a Viernes.
    \item El horario de atención a estudiantes en oficina (Office Hours) será los días Martes de 10:00 a 11:00am durante el transcurso del semestre académico. La atención se realizará por orden de llegada (FIFO).
\end{itemize}

\vspace{4em}

% --- Signatures ---
\begin{minipage}[t]{0.45\textwidth}

    \input{signature}

    \vspace{-17mm}
    \noindent \rule{\textwidth}{0.4pt} \\ % Línea de firma
    \textbf{Provider Signature} \\
    Date: \makebox[1.5in]{\textbf{\today}}
\end{minipage}
\hfill
\begin{minipage}[t]{0.45\textwidth}
    \vspace{-6mm}
    \noindent \rule{\textwidth}{0.4pt} \\
    \textbf{Client Signature} \\
    Date: \makebox[1.5in]{\hrulefill}
\end{minipage}

\end{document}
